\section*{Bab \ref{ch:b3:particle-physics} --- Fisika Partikel}

\begin{solution}{exo:b3:particle-physics:1}
(a) uud: $\tfrac23 + \tfrac23 - \tfrac13 = +1$; udd: $\tfrac23 -
\tfrac13 - \tfrac13 = 0$. (b) $\tfrac23 + \tfrac13 = +1$ (yakni
anti-d yang membawa $+\tfrac13$). (c) $\pi^- =
\bar{\text{u}}\text{d}$; $\Delta^{++} =$ uuu ($3\times\tfrac23 =
+2$ --- yakni keadaan yang statistikanya membantu memaksa penciptaan
warna). (d) Medan gluon yang terentang menyimpan energi
$\propto$ panjangnya; jauh sebelum sebutir \omterm{def:b3:particle-physics:zoo}{kuark} bebas, energi tersimpannya
melampaui $2m_{\text{q}}c^2$ dan talinya patah menjadi pasangan
yang baru: Anda menarik keluar sebutir \omterm{def:b3:particle-physics:zoo}{meson}, tak pernah \omterm{def:b3:particle-physics:zoo}{kuark} telanjang.
\end{solution}

\begin{solution}{exo:b3:particle-physics:2}
(a) \omterm{def:b3:particle-physics:zoo}{Lepton}: e$^-$, $\nu_\mu$, $\mu^-$; \omterm{def:b3:particle-physics:zoo}{barion}: p, n; \omterm{def:b3:particle-physics:zoo}{meson}:
$\pi^+$. (b) Hanya \omterm{def:b3:particle-physics:zoo}{hadronnya}: $\pi^+$, p, n. (c) Yang mantap sendirian:
e$^-$, p, dan neutrinonya (neutron bebas hidup lima belas
menit; $\mu$ dan $\pi$ mati dalam hitungan mikrodetik dan kurang). (d) Yang
netral: $\nu_\mu$ dan n.
\end{solution}

\begin{solution}{exo:b3:particle-physics:3}
(a) Diizinkan --- yakni peluruhan $\beta^-$ yang baku (neutron bebas
melakukannya, $T_{1/2} \approx \qty{10}{min}$). (b) Dilarang oleh energinya:
hasilnya lebih berat daripada proton bebasnya (di dalam sebuah inti, \omterm{prop:b3:nuclear-physics:binding}{energi
ikatnya} dapat membayar, dan peluruhan $\beta^+$ yang terikat memang terjadi). (c)
Dilarang: \omterm{prop:b3:particle-physics:conservation}{bilangan barionnya} $1 \to 0$ (dan \omterm{prop:b3:particle-physics:conservation}{bilangan leptonnya}
$0 \to -1$). (d) Dilarang oleh \omterm{prop:b3:particle-physics:conservation}{bilangan lepton} \emph{keluarganya}:
kemuonannya akan lenyap dan keelektronannya muncul --- yang tak pernah
teramati, dan diburu justru karena alasan itu.
\end{solution}

\begin{solution}{exo:b3:particle-physics:4}
(a) $6\times0.938 \approx \qty{5.6}{GeV}$. (b) Energi kinetiknya
$m_{\text{p}}c^2 \approx \qty{0.94}{GeV}$ tiap berkasnya
($\sqrt s = 4m_{\text{p}}c^2$ secara berhadapan). (c) Kira-kira $3:1$:
kekekalan momentum menghukum sebagian besar energi sasaran tetapnya menjadi
gerak maju yang tak berguna bagi puingnya --- berkas yang berhadapan
membelanjakan segalanya. (d) Mesin tahun 1955 dirancang tepat melewati
ambang \qty{5.6}{GeV}: anggaran pemercepat ditulis dalam
massa tak berubah, dan antiprotonnya tiba sesuai jadwal.
\end{solution}

\begin{solution}{exo:b3:particle-physics:5}
(a) $c\tau = \qty{660}{m}$. (b) Jangkauan terulur sebesar
${\sim}\qty{15}{km}$ memerlukan $\gamma \sim 23$: $E = \gamma
m_\mu c^2 \sim \qty{2.4}{GeV}$ --- persis skala kosmiknya. (c)
Dalam kerangkanya sendiri muonnya hidup \qty{2.2}{\micro s} biasanya
sementara atmosfernya melesat lewat, terkerut menjadi
${\sim}\qty{650}{m}$: kelangsungan hidup yang sama, buku besar yang lain. (d) \omterm{prop:b3:particle-physics:conservation}{Bilangan
barionnya}: tak ada \omterm{def:b3:particle-physics:zoo}{barion} yang cukup ringan, dan muonnya adalah sebutir
\omterm{def:b3:particle-physics:zoo}{lepton}; peluruhannya harus mengubah rasa, yang hanya dilakukan gaya lemah
--- sehingga muncullah mikrodetik yang anggun dan bukan
yoktodetik gaya kuatnya.
\end{solution}

\begin{solution}{exo:b3:particle-physics:6}
(a) $\text{d} \to \text{u} + \text{e}^- + \bar\nu_{\text{e}}$:
muatannya $-\tfrac13 = \tfrac23 - 1 + 0$; \omterm{prop:b3:particle-physics:conservation}{bilangan barionnya}
$\tfrac13 = \tfrac13$; bilangan keluarga elektronnya $0 = 1 - 1$. (b) Peluruhan
dua benda menetapkan energi elektronnya; \emph{kesinambungan} yang
teramati berarti entah kekekalan energinya gagal
(Bohr sempat mempertimbangkannya) atau benda ketiga yang tak terlihat berbagi
anggarannya --- yakni neutrino Pauli, yang terdeteksi dua puluh lima tahun kemudian.
(c) $\hbar/m_{\text{W}}c = 197/80400 \approx
\qty{2.5e-3}{fm}$ --- yakni seperseribu sebutir proton. (d) Pada energi
rendah interaksinya harus menjembatani celah yang dibuat sangat pendek oleh
perantaranya yang berat, sehingga amplitudonya mungil; di dekat
\qty{100}{GeV} W-nya terjangkau dan gaya ``lemah'' itu menunjukkan
kekuatan elektrolemahnya yang sejati.
\end{solution}

\begin{solution}{exo:b3:particle-physics:7}
(a) Gaya kuatnya mengekalkan keanehan, sehingga ia hanya dapat membuat
$+1$ dan $-1$ bersama-sama: yakni sebutir kaon dengan hiperonnya. (b) Tiga belas
orde besar terlalu lambat bagi gaya kuatnya: gaya
lemahnya, satu-satunya pematah rasa, menandatangani surat kematiannya.
(c) $10^{13}$: yakni satu detik berbanding tiga ratus ribu tahun.
(d) Kelahiran yang berlimpah berkata ``kuat''; kematian yang enggan berkata ``bukan
kuat'' --- hanya sebuah bilangan kekal baru, yang dibuat berpasangan oleh gaya
kuatnya dan dipatahkan oleh yang lemah, yang menyerasikan keduanya: yakni
pembukuannya yang \emph{merupakan} penemuannya.
\end{solution}

\begin{solution}{exo:b3:particle-physics:8}
(a) Dalam kerangka diam pasangannya momentum totalnya nol: satu
foton tak pernah dapat bermomentum nol, dua yang saling membelakangi dapat ---
kekekalan menuliskan geometri PET. (b) $E = mc^2 =
0.002\times(3\times10^8)^2 = \qty{1.8e14}{J}$: dua gram
pemusnahan mengalahkan satu kilogram uranium yang terfisi. (c) Dari
\omterm{def:b3:nuclear-physics:nuclide}{isotop} $\beta^+$ buatan siklotron --- yakni $^{18}$F, yang dibangun ke dalam
glukosa: kimia tubuhnya sendiri menempatkan sumber positronnya dalam
sel yang paling lapar. (d) Bahan bakar yang berongkos sejuta kali
hasilnya, tanpa tambang alami: antimateri adalah sebuah kalibrasi
pendeteksi dan mesin sang pendongeng, bukan sumber energi.
\end{solution}

\begin{solution}{exo:b3:particle-physics:9}
(a) Potensial Higgsnya adalah sumur ganda (yakni topi Meksiko): yakni
titik setangkupnya $\phi = 0$ adalah sebuah puncak, dan vakumnya menggelinding
ke tepinya --- keadaan dasar alam semesta mematahkan kesetangkupannya
sendiri, persis seperti magnet yang memilih utara. (b) Pilihannya
hanya mematahkan sebagian kesetangkupannya: gabungan yang selamat
adalah elektromagnetisme, sehingga pembawanya --- yakni fotonnya --- menyimpan
massa nol, sedangkan W dan Z, pembawa arah yang patah,
memakan kekakuan yang bersesuaian itu sebagai massa. (c) Getaran radial
\omterm{def:b3:phase-transitions:order}{parameter keteraturannya} --- dalam magnetnya, yakni
fluktuasi $|m|$ di sekitar nilai kesetimbangannya di dekat
$T_{\text{c}}$. (d) Dalam $10^{-11}$ detik pertama alam
semesta (\cref{ch:b3:astrophysics}): di atas suhu
elektrolemahnya tiap partikelnya tak bermassa dan kedua gayanya
adalah satu.
\end{solution}

\begin{solution}{exo:b3:particle-physics:10}
(a) Tiap \qty{26.7}{MeV} $= \qty{4.28e-12}{J}$ menyerahkan dua
neutrino: fluksnya $= 2\times1360/4.28\times10^{-12} \approx
\qty{6.4e14}{m^{-2}.s^{-1}}$. (b) Menembus sebuah kuku ibu jari:
$6\times10^{10}$ tiap detiknya. Lintasan seumur hidup menembus tubuh Anda
yang ${\sim}\qty{0.03}{m^2}$ dan setebal \qty{0.3}{m} selama
${\sim}\qty{2.5e9}{s}$: kira-kira $5\times10^{22}$, dikalikan
$10^{-22}\,\unit{m^{-1}}\times\qty{0.3}{m}$: berorde satu ---
satu neutrino matahari akan menyentuh Anda seumur hidup Anda. (c)
$\nu_{\text{e}}$ yang hilang telah berayun menjadi $\nu_\mu$ dan
$\nu_\tau$ di perjalanan; ayunannya menuntut selisih massa, sehingga
neutrinonya bermassa --- yakni retakan laboratorium pertama pada Model
Baku yang berneutrino tak bermassa. (d) Gunungnya menyaring
segala hal \emph{lainnya}: hanya neutrino yang lewat, sehingga berkilo-kilometer
batuan mengubah langit yang paling berisik menjadi laboratorium yang paling sunyi.
\end{solution}

\begin{solution}{exo:b3:particle-physics:11}
(a) $\gamma = 6.8\times10^{12}/9.38\times10^{8} \approx 7250$:
$1 - v/c \approx 1/2\gamma^2 \approx 10^{-8}$ --- yakni tiga meter
tiap detik kurang dari cahaya. (b) $\lambda \approx hc/E \approx
\qty{3e-5}{fm}$: yakni sepertiga puluh ribu jejari proton --- wilayah
\omterm{def:b3:particle-physics:zoo}{kuark}. (c) $13.6\,\text{TeV}/938\,\text{MeV} \approx
14500$ massa proton --- seandainya seluruhnya dapat diubah. (d) Pada
energi tetap $\gamma_{\text{e}}/\gamma_{\text{p}} =
m_{\text{p}}/m_{\text{e}} \approx 1836$, sehingga rugi elektronnya
$1836^4 \sim 10^{13}$ kali lebih buruk: proton untuk cincin bertenaga kasar,
elektron (yang bersih dan bertitik) untuk ketelitian pada energi yang lebih rendah.
\end{solution}

\begin{solution}{exo:b3:particle-physics:12}
(a) $3\times10^{20}\times1.6\times10^{-19} \approx \qty{48}{J}$:
yakni sebuah bola tenis yang diservis keras --- dalam satu proton. (b) $\gamma \approx
3\times10^{11}$; $10^{5}$ tahun cahaya Galaksinya terkerut menjadi
${\sim}10$ \emph{detik} cahaya: ia menyeberangi Bima Sakti, menurut
jamnya sendiri, dalam hitungan detik. (c) Dalam kerangka protonnya, gelombang
mikro CMB yang jinak tiba tergeser biru menjadi sinar $\gamma$ yang cukup
energetik untuk memukul pion darinya: di atas ambangnya alam semesta
sendiri menjadi penyerap, dan proton semacam itu pastilah setempat. (d) Satu
partikel semacam itu tiap kilometer persegi tiap abadnya, dari jenis yang tak diketahui,
yang tak dibidikkan siapa pun: $10^{9}$ tumbukan tiap detiknya di LHC,
dari berkas yang dikenal pada energi yang dikenal, membeli statistika dan
kendali yang tak dapat ditandingi hadiah kosmik mana pun.
\end{solution}

\begin{solution}{pb:b3:particle-physics:1}
\textbf{1.} Sebutir proton primer (sering \unit{GeV}--\unit{TeV})
menghantam sebuah inti jauh di atas (lewat gaya kuat), lalu menyemburkan pion;
pion bermuatannya meluruh lewat gaya lemah menjadi muon dan neutrino;
muonnya, yang mengionkan hanya secara lembut (secara elektromagnetik), berpacu menuju
tanah.
\textbf{2.} $c\tau_\pi \approx \qty{8}{m}$: bahkan $\gamma$ yang besar
hanya membeli ratusan meter bagi pionnya, dan mereka juga mati lewat
\omterm{rem:b3:particle-physics:forces}{interaksi kuat} dalam udaranya; muonnya, yang tak merasakan gaya kuat dan
hidup $85\times$ lebih lama, adalah yang menempuh perjalanan.
\textbf{3.} $400\,\unit{cm^2}\times1\,\text{min}^{-1} = 400$
tiap menitnya --- yakni kira-kira 7 tiap detiknya.
\textbf{4.} Kilatan derau dan keradioaktifan setempat menyalakan satu
dayung secara acak; menyalakan \emph{keduanya} dalam \qty{100}{ns}
menuntut satu partikel tunggal yang melintasi pasangannya --- yakni geometri sebagai
sebuah tapis.
\textbf{5.} $R_{\text{acc}} = R_1R_2\Delta t = 100\times100
\times10^{-7} = \qty{e-3}{s^{-1}}$ --- yakni satu cacahan palsu tiap
seperempat jam berbanding seratus muon tiap menitnya: dapat diabaikan.
\textbf{6.} Hanya muon dalam kerucut yang menembus kedua
dayungnya (yang nyaris tegak, dalam ${\sim}30$--$40^\circ$) yang
diterima: pasangannya mengukur sebuah arah, bukan sekadar laju ---
dan itulah yang menjadikannya sebuah \emph{teleskop}.
\textbf{7.} Muon yang miring menyeberangi lebih banyak atmosfer: lebih banyak rugi
energi dan lebih banyak \omterm{prop:b3:relativistic-kinematics:dilation}{waktu diri} dalam penerbangannya, sehingga lebih sedikit yang selamat ---
kira-kira hukum $\cos^2\theta$ yang dijejak percobaan miring Anda.
\textbf{8.} $c\tau = \qty{660}{m}$: muon klasik dari
\qty{15}{km} selamat dengan peluang
$\eu^{-15000/660} \approx 10^{-10}$ --- lotengnya akan mencacah
satu muon tiap dasawarsa.
\textbf{9.} $\gamma = 4000/106 \approx 38$.
\textbf{10.} Umur terulurnya $\gamma\tau \approx
\qty{84}{\micro s}$: jangkauannya ${\sim}\qty{25}{km}$ --- laju
lotengnya persis seperti yang diramalkan kenisbian.
\textbf{11.} Dalam kerangka muonnya umurnya adalah \qty{2.2}{\micro s}
biasa tetapi atmosfernya terkerut menjadi
$15000/38 \approx \qty{400}{m}$: terseberangi dengan waktu bersisa.
\textbf{12.} $t_{\text{diri}} \approx 15000/(38\times3\times
10^8) = \qty{1.3}{\micro s}$: kelangsungan hidupnya $\eu^{-1.3/2.2} \approx
0.55$ --- separuhnya tiba, dan pencacah Anda adalah sensus mereka.
\textbf{13.} Yakni percobaan muon klasik gunung berbanding aras laut
(Rossi--Hall, dan Frisch--Smith yang terfilmkan): peluruhan
klasiknya meramalkan pada dasarnya nol cacahan; ratusan cacahan tiap menit
Anda adalah kenisbian khusus, yang dibuktikan di atas \omterm{def:b3:electrons-in-solids:classes}{penyekat} atapnya.
\textbf{14.} Sebutir \omterm{def:b3:particle-physics:zoo}{lepton} bermuatan dari keturunan kedua:
bermuatan $-e$, berspin $\tfrac12$, bermassa \qty{106}{MeV} --- yakni
kembaran berat elektron.
\textbf{15.} Muatannya: $-1 = -1 + 0 + 0$; keluarga elektronnya:
$0 = 1 + (-1) + 0$; keluarga muonnya: $1 = 0 + 0 + 1$. Semua bukunya
seimbang.
\textbf{16.} Tiga benda berbagi anggarannya dalam segala perbandingan,
sehingga spektrum elektronnya bersifat sinambung; ia berakhir ketika kedua
neutrinonya melenting bersama melawan elektronnya: yakni kira-kira
$m_\mu c^2/2 \approx \qty{53}{MeV}$.
\textbf{17.} Ia akan mematahkan bilangan keluarga \omterm{def:b3:particle-physics:zoo}{lepton} --- yakni satu
buku yang diseimbangkan Model Baku tanpa tahu mengapa. Menemukannya
akan menjadi sepupu sektor bermuatan bagi ayunan neutrino:
yakni fisika baru yang terjamin, dan karena itulah percobaan mengejarnya sampai
satu bagian dalam $10^{13}$.
\textbf{18.} Masuknya muonnya membuat satu kilatan; beberapa mikrodetik
kemudian elektron peluruhannya membuat kilatan kedua dalam dayung yang sama.
Histogram jedanya adalah sebuah eksponensial yang kemiringannya \emph{adalah}
$\tau = \qty{2.2}{\micro s}$: yakni sebuah umur yang terukur di atas rak.
\textbf{19.} Semua rumus Bohr berskala dengan massanya: orbit
hidrogen muoniknya menyusut sebesar $207$ menjadi ${\sim}\qty{250}{fm}$ ---
begitu dalam di dalam awan elektronnya sehingga muonnya merasakan
ukuran protonnya: atom muonik adalah penggaris proton.
\textbf{20.} $10^{15}/10^{6} = \qty{e9}{eV}$: yakni kira-kira satu GeV tiap
partikelnya --- muon loteng Anda adalah warga hujan yang rata-rata.
\textbf{21.} Muon relativistik mendahului cahaya dalam air lalu berpijar
biru Cherenkov ke dalam pengganda foton tangkinya. Tapak sebuah hujan
$10^{20}\,\unit{eV}$ terbentang berkilo-kilometer persegi, sehingga tangki yang
jarang mencupliknya seperti tetes hujan dalam
ember.
\textbf{22.} Sebuah partikel menaruh energinya dalam medium yang
mengubahnya menjadi cahaya, sebuah pengganda foton mengubah cahayanya menjadi
elektron, dan elektronikanya mencacahnya: tiap kalorimeter dari
loteng sampai penumbuk adalah kalimat ini.
\textbf{23.} Jejari lingkarannya memberikan momentumnya, arah putarnya memberikan
tanda muatannya; foton, neutron dan neutrinonya ---
yang netral --- tak meninggalkan jejak, hanya utangnya dalam
pembukuannya.
\textbf{24.} Alam mempercepat proton tunggal sepuluh juta
kali melampaui LHC --- yang merendahkan hati --- tetapi pada satu tiap kilometer
persegi tiap abadnya: mesin kita menukar energi puncaknya dengan
$10^{9}$ tumbukan terkendali tiap detiknya.
\textbf{25.} Lahir dari sebutir proton sepuluh kilometer di atas; diantarkan
hidup-hidup oleh \omterm{prop:b3:relativistic-kinematics:dilation}{pemuluran waktu} ($\gamma \approx 38$, separuhnya selamat);
dikenali oleh dua kilatan yang berjarak seratus nanodetik; dan
dijelaskan, seluruhnya, oleh sebuah tabel berisi dua belas kursi dan empat
gaya.
\end{solution}
